% EABC_Monopol_Artikel.tex
% Standalone-Preprint: EABC-Monopol und spektrale Riemann-Nullstellen-Analyse
% Kompilierung: pdflatex EABC_Monopol_Artikel.tex

\documentclass[11pt,a4paper]{article}

\usepackage[utf8]{inputenc}
\usepackage[T1]{fontenc}
\usepackage[ngerman]{babel}
\usepackage{amsmath,amssymb}
\usepackage{booktabs}
\usepackage{xcolor}
\usepackage{microtype}
\usepackage{hyperref}

\hypersetup{
  colorlinks=true,
  linkcolor=blue!60!black,
  urlcolor=blue!60!black,
  citecolor=blue!60!black,
}

% --- Ergebnismakros N=150 (Referenzlauf EABC_Gross_Analyse.py) ---
\input{monopol_preprint_results.txt}

% --- Modellkonstanten (Postulat, nicht aus E8 abgeleitet) ---
\newcommand{\monopolQtot}{1796}
\newcommand{\monopolFactor}{1036.92} % 1796/\sqrt{3}

% --- Konvergenz Teilstichproben (gleicher N=150-Lauf) ---
\newcommand{\convFiftyVar}{0.172806}
\newcommand{\convFiftyDelta}{0.068616}
\newcommand{\convFiftyMean}{0.992296}
\newcommand{\convHundredVar}{0.155540}
\newcommand{\convHundredDelta}{0.051350}
\newcommand{\convHundredMean}{0.995738}

% --- Fehlerhafte Hauptterm-Pipeline (Diagnose, t in [5000,5140]) ---
\newcommand{\badMeanFixed}{2.9586}
\newcommand{\badVarFixed}{7.33}
\newcommand{\badMeanAdapt}{2.9962}
\newcommand{\badVarAdapt}{7.88}
\newcommand{\badCountInterval}{52}
\newcommand{\goodCountInterval}{150}

% --- Ergebnismakros N=500 (Referenzlauf Monopol Test.py / n500_reference) ---
\newcommand{\bulkFiveHundredN}{500}
\newcommand{\bulkFiveHundredMean}{0.999869}
\newcommand{\bulkFiveHundredVar}{0.157876}
\newcommand{\bulkFiveHundredDelta}{0.053686}
\newcommand{\bulkFiveHundredPm}{0.032292}
\newcommand{\bulkFiveHundredPressureSum}{16.145954}
\newcommand{\bulkFiveHundredTmax}{5466.4183}
\newcommand{\bulkFiveHundredRuntime}{106.6}

% --- Ergebnismakros N=1000 (bulk_results.json / EABC_Bulk_Evaluation.py) ---
\newcommand{\bulkThousandN}{1000}
\newcommand{\bulkThousandMean}{1.004083}
\newcommand{\bulkThousandVar}{0.161290}
\newcommand{\bulkThousandDelta}{0.057099}
\newcommand{\bulkThousandPm}{0.018441}
\newcommand{\bulkThousandTmax}{5931.5701}
\newcommand{\bulkThousandRuntime}{606.9}
\newcommand{\bulkGUE}{0.104192}

% Block-Stabilität N=1000 (je 250 Nullstellen)
\newcommand{\blkOneMean}{1.007297}
\newcommand{\blkOneVar}{0.169773}
\newcommand{\blkOneDelta}{0.065581}
\newcommand{\blkOnePm}{0.130777}
\newcommand{\blkTwoMean}{1.001638}
\newcommand{\blkTwoVar}{0.156151}
\newcommand{\blkTwoDelta}{0.051960}
\newcommand{\blkTwoPm}{-0.057014}
\newcommand{\blkThreeMean}{1.008279}
\newcommand{\blkThreeVar}{0.165494}
\newcommand{\blkThreeDelta}{0.061302}
\newcommand{\blkThreePm}{-0.119749}
\newcommand{\blkFourMean}{0.999098}
\newcommand{\blkFourVar}{0.153654}
\newcommand{\blkFourDelta}{0.049463}
\newcommand{\blkFourPm}{0.124070}

% --- t'Hooft-Monopolmasse (Referenzlauf hooft.py) ---
\newcommand{\hooftAlphaInv}{137.035999177}
\newcommand{\hooftMX}{2.0\times 10^{16}\,\mathrm{GeV}}
\newcommand{\hooftMPlanck}{1.2209\times 10^{19}\,\mathrm{GeV}}
\newcommand{\hooftMMPlanck}{0.224484}
\newcommand{\hooftAlphaG}{0.050393}
\newcommand{\hooftInvAlphaG}{19.84}

\title{EABC-Magnetmonopol und spektrale Abweichungen\\
  der Riemann-Nullstellen im hochenergetischen Fernfeld}
\author{Thomas Hoffbauer\\[0.5em]
  \small\textit{Preprint-Entwurf; numerische Begleitdokumentation zum EABC-Modell}}
\date{\today}

\begin{document}
\maketitle

\begin{abstract}
  Im Rahmen des EABC-Frameworks (Entanglement-ABC-Klassifikation quaternionischer
  Zustände) wird ein \emph{Modell} magnetischer Monopole in zwei komplementären
  Bildern diskutiert: ein globales Punkt-Monopol mit Postulat
  $Q_{\mathrm{tot}}=1796$ sowie diskrete topologische Defekte im Hurwitz-Quaternionen-Gitter
  (\texttt{MagMo.py}). Diese Größen sind \emph{nicht} aus einer E$_8$-Konstruktion
  abgeleitet und stellen Modellparameter dar.

  Zur empirischen Prüfung wird das Spektrum aufeinanderfolgender nichttrivialer
  Riemann-Nullstellen ab $t_0=5000$ ausgewertet. Nach methodischer Korrektur
  (vollständige \texttt{mpmath.siegelz}-Pipeline statt unvollständigem Siegel-Hauptterm)
  bleibt der mittlere normierte Abstand Weyl-konform ($\bar{s}\approx 1$); eine
  früher diskutierte ``Drittel-Regel'' der Nullstellendichte wird verworfen.
  Die gesicherte Abweichung vom GUE-Referenzensemble zeigt sich im \emph{zweiten Moment}:
  $\Delta\sigma^2 \approx +0{,}053$--$+0{,}057$ für Stichproben $N=150$, $500$ und $1000$.
  Der normierte spektrale Monopol-Druck $P_m$ oszilliert stark und ist blockweise instabil.

  Der Text erhebt keinen Anspruch auf experimentellen Nachweis eines magnetischen
  Monopols im Labor; er dokumentiert numerische Korrelationen innerhalb eines
  offenen theoretischen Modells.
\end{abstract}

%=============================================================================
\section{Einleitung}
\label{sec:einleitung}
%=============================================================================

Das EABC-Framework ordnet arithmetische und spektrale Strukturen über eine
quaternionische Vier-Kanal-Klassifikation ($E,A,B,C$) und diskrete Gittermodelle.
Im Kontext der Riemannschen Zeta-Funktion $\zeta(s)$ auf der kritischen Linie
$\Re(s)=\tfrac12$ stellt sich die Frage, ob zusätzliche ``topologische'' Freiheitsgrade
--- formal als magnetische Monopole modelliert --- im Nullstellen-Spektrum messbare
Spuren hinterlassen.

\paragraph{Monopol-Hypothese.}
Wir unterscheiden zwei \emph{Modellbilder}, die derzeit \emph{nicht} in einer
einheitlichen Feldtheorie verknüpft sind:
\begin{enumerate}
  \item ein \textbf{globales Punkt-Monopol} mit effektiver Gesamtladung
        $Q_{\mathrm{tot}}=1796$ als Skalierungsparameter der spektralen Kopplung;
  \item \textbf{diskrete Hurwitz-Defekte} im Quaternionen-Primgitter, identifiziert
        durch verschwindenden Skalarteil reiner Imaginär-Primzahlen
        (\texttt{MagMo.py}).
\end{enumerate}
Beide Bilder dienen der \emph{Hypothesenbildung}; keines davon folgt rigoros aus
E$_8$ oder aus der analytischen Zahlentheorie der $\zeta$-Funktion.

\paragraph{Numerische Strategie.}
Als empirischer Zugang werden lückenlose Ketten von Nullstellen $\{\gamma_n\}$
im Fernfeld $t\gtrsim 5000$ berechnet, über das Weyl-Gesetz entfaltet und mit
der Wigner-Dyson-GUE-Referenz ($\sigma^2_{\mathrm{GUE}}\approx 0{,}104$) verglichen.
Implementierungsdetails finden sich in \texttt{EABC\_Gross\_Analyse.py} ($N=150$),
\texttt{Monopol Test.py} ($N=500$) und \texttt{EABC\_Bulk\_Evaluation.py} ($N=1000$).

%=============================================================================
\section{Eigenschaften des EABC-Monopols}
\label{sec:monopol-eigenschaften}
%=============================================================================

\subsection{Globales Punkt-Monopol ($Q_{\mathrm{tot}}=1796$)}

Im spektralen Testmodell wird ein dominantes, nicht kompensiertes Punkt-Monopol
durch die \emph{Modellkonstante}
\begin{equation}
  Q_{\mathrm{tot}} = 1796
  \label{eq:qtot}
\end{equation}
charakterisiert. Dieser Wert ist ein \textbf{Postulat} des EABC-Monopolmodells;
eine Herleitung aus der Wurzelsystem-Dimension von E$_8$ oder aus einer
Gauge-Theorie-Konstruktion liegt nicht vor.

Daraus folgt der dimensionslose Skalierungsfaktor
\begin{equation}
  f_{\mathrm{mon}} = \frac{Q_{\mathrm{tot}}}{\sqrt{3}} \approx \monopolFactor,
  \label{eq:fmon}
\end{equation}
der in den numerischen Skripten als ``asymptotische Feldkonstante'' für die
Kopplung an die Nullstellen $\gamma_n$ verwendet wird. Physikalisch wäre ein
magnetischer Monopol in einer kompakten Eichtheorie durch Dirac-Quantisierung
diskret; hier wird \emph{kein} solcher Nachweis behauptet --- $Q_{\mathrm{tot}}$
parametrisiert lediglich die Modellamplitude.

\subsection{Hurwitz-Defekte im Quaternionen-Gitter (\texttt{MagMo.py})}

Komplementär betrachtet \texttt{MagMo.py} Hurwitz-Quaternionen-Primzahlen bis
einer Normschranke und isoliert ``Monopol-Kandidaten'' als Gitterpunkte mit
verschwindendem Skalarteil $a_0=0$ (rein imaginäre Hurwitz-Primzahlen).
Die akkumulierte ``Feld-Divergenz'' an einem Gitterknoten $(i,j,k)$ zählt
lokale topologische Defekte; dieses Bild ist \emph{diskret} und \emph{ortsaufgelöst},
im Gegensatz zum globalen $Q_{\mathrm{tot}}$-Postulat.

Die beiden Beschreibungen --- globales $Q=1796$ und verteilte Hurwitz-Defekte ---
sind derzeit \emph{nicht} vereinheitlicht. Eine offene Aufgabe ist die Abbildung
der Gitter-Defektdichte auf einen effektiven spektralen Kopplungsterm
(Abschnitt~\ref{sec:offene-fragen}).

\subsection{Spektraler Monopol-Druck}

Als aggregierte, stichprobenabhängige Größe definieren wir die \emph{Summe}
\begin{equation}
  S_N = \sum_{n=1}^{N} \sin\!\Bigl(\gamma_n \,\frac{\log f_{\mathrm{mon}}}{150}\Bigr)
  \label{eq:pressure-sum}
\end{equation}
und den \emph{normierten} Druck
\begin{equation}
  P_m = \frac{S_N}{N}.
  \label{eq:pressure-pm}
\end{equation}
Wegen der Oszillation einzelner Terme ist $P_m$ \emph{nicht} stabil gegenüber
kleinen Stichproben und wechselt das Vorzeichen zwischen Datenblöcken
(Tabelle~\ref{tab:block-stabilitaet}). Für Vergleiche ist stets $N$ anzugeben.

%=============================================================================
\section{t'Hooft-Monopolmasse und EABC-Gravitationsstrukturkonstante}
\label{sec:hooft-monopolmasse}
%=============================================================================

Unabhängig vom spektralen EABC-Testmodell liefert die semiklassische
t'Hooft--Polyakov-Konstruktion~\cite{thooft,polyakov} in einer
SU(2)-Eichtheorie mit spontan gebrochener Symmetrie magnetische Monopole
als topologisch stabile Solitonen. Ihre Masse skaliert mit der
Symmetriebrechungsskala $M_X$ und der feinstrukturierten Kopplungskonstante
$\alpha = e^2/(4\pi\varepsilon_0\hbar c)$:
\begin{equation}
  M_M \sim \frac{M_X}{\alpha}.
  \label{eq:hooft-mass}
\end{equation}
Diese Relation ist \emph{nicht} aus dem Postulat $Q_{\mathrm{tot}}=1796$
herleitbar; sie dient hier als externer Größenordnungsvergleich für
physikalische Monopolmassen.

\paragraph{Reproduzierbare Berechnung (\texttt{hooft.py}).}
Das Begleitskript \texttt{hooft.py} wertet \eqref{eq:hooft-mass} mit den
Eingabewerten $\alpha^{-1}=\hooftAlphaInv$, $M_X=\hooftMX$ und
$M_{\mathrm{Planck}}=\hooftMPlanck$ aus und bildet das dimensionslose Verhältnis
$M_M/M_{\mathrm{Planck}}$ sowie die EABC-Gravitationsstrukturkonstante
\begin{equation}
  \alpha_G^{\mathrm{EABC}} \equiv \Bigl(\frac{M_M}{M_{\mathrm{Planck}}}\Bigr)^2.
  \label{eq:alpha-g-eabc}
\end{equation}
Tabelle~\ref{tab:hooft-ergebnisse} fasst die numerischen Resultate zusammen.

\begin{table}[htbp]
\centering
\caption{t'Hooft-Monopolmasse und abgeleitete EABC-Größen
(\texttt{hooft.py}, $36$-stellige Dezimalarithmetik).}
\label{tab:hooft-ergebnisse}
\begin{tabular}{ll}
\toprule
\textbf{Größe} & \textbf{Wert} \\
\midrule
$\alpha^{-1}$ & $\hooftAlphaInv$ \\
$M_X$ & $\hooftMX$ \\
$M_{\mathrm{Planck}}$ & $\hooftMPlanck$ \\
$M_M/M_{\mathrm{Planck}}$ & $\hooftMMPlanck$ \\
$\alpha_G^{\mathrm{EABC}} = (M_M/M_{\mathrm{Planck}})^2$ & $\hooftAlphaG$ \\
$1/\alpha_G^{\mathrm{EABC}}$ & $\hooftInvAlphaG$ \\
\bottomrule
\end{tabular}
\end{table}

\paragraph{Vorsichtige Einordnung im EABC-Framework.}
Die Größe $\alpha_G^{\mathrm{EABC}}\approx 0{,}050$ bzw.\ ihr Kehrwert
$1/\alpha_G^{\mathrm{EABC}}\approx 19{,}8$ charakterisiert im EABC-Bild eine
dimensionslose \emph{Gravitationsstrukturkonstante} auf Planck-Skala.
Sie entsteht allein aus der t'Hooft-Massenskala und $M_{\mathrm{Planck}}$;
eine numerische Koinzidenz oder Identifikation mit dem spektralen Postulat
$Q_{\mathrm{tot}}=1796$ (Abschnitt~\ref{sec:monopol-eigenschaften}) wird
\emph{nicht} behauptet. Beide Konstanten gehören zu verschiedenen
Modellschichten --- $Q_{\mathrm{tot}}$ parametrisiert die Kopplung an
Riemann-Nullstellen, während $M_M$ eine Feldtheorie-Massenskala beschreibt.
Eine Vereinheitlichung bleibt Gegenstand offener Fragen
(Abschnitt~\ref{sec:offene-fragen}).

%=============================================================================
\section{Methodische Richtigstellung}
\label{sec:methodische-richtigstellung}
%=============================================================================

In der numerischen Verifikation im hochenergetischen Bereich ($t\approx 5000$)
zeigt sich eine methodische Divergenz, die vor jeder spektralen Auswertung
geklärt werden muss.
Eine Nullstellensuche, die nur den asymptotischen Siegel-\emph{Hauptterm} nutzt
und den Restterm $R(t)$ vernachlässigt, erzeugt systematische Lücken in der Stichprobe.

\begin{table}[htbp]
\centering
\caption{Vergleich numerischer Approximationsmethoden im Intervall $t \in [5000,\,5140]$.
Die ersten beiden Zeilen verwenden den unvollständigen Hauptterm; die dritte Zeile die
vollständige Implementierung \texttt{mpmath.siegelz} mit adaptiver Schrittweite
und \texttt{brentq}/Illinois-Verfeinerung.}
\label{tab:siegel_vergleich}
\begin{tabular}{lccc}
\toprule
\textbf{Methode} &
\textbf{Mittlerer norm.\ Abstand} &
\textbf{Varianz $\sigma^2$} &
\textbf{Gefundene Nullstellen} \\
\midrule
Hauptterm, Schrittweite $0{,}5$ &
$\badMeanFixed$ & $\badVarFixed$ & $\sim\badCountInterval$ \\
Hauptterm, adaptive Schrittweite &
$\badMeanAdapt$ & $\badVarAdapt$ & $\sim\badCountInterval$ \\
\texttt{mpmath.siegelz}, adaptiv &
$\convFiftyMean$ & $\convFiftyVar$ & $\sim\goodCountInterval$ \\
\bottomrule
\end{tabular}
\end{table}

Tabelle~\ref{tab:siegel_vergleich} zeigt: Eine rein adaptive Schrittweite behebt das
Problem nicht, solange $R(t)$ fehlt --- es werden dieselben $\sim\badCountInterval$ Nullstellen
gefunden wie mit fester Schrittweite $0{,}5$.
Der scheinbare Faktor $\sim 3$ im mittleren normierten Abstand ist ein
Aliasing-Effekt der Unterabtastung, kein spektraler Beitrag des Monopols.

Erst unter Einbeziehung der vollständigen Riemann-Siegel-Asymptotik stellt sich
Weyl-konformes Verhalten mit mittlerem Abstand $\approx 1$ ein.
Eine Modifikation der globalen Zählfunktion der Form
$N_{\mathrm{EABC}}(T) \approx \tfrac{1}{3}\,N_{\mathrm{standard}}(T)$
--- die sogenannte ``Drittel-Regel'' oder jede daraus abgeleitete
``Drittel-Symmetrie'' der Nullstellendichte ---
ist ohne unabhängige Evidenz \textbf{nicht haltbar und wird verworfen}.

\paragraph{Korrigierte Pipeline.}
Alle im Folgenden berichteten Ergebnisse verwenden \texttt{mpmath.siegelz} mit
konservativer Schrittweite (ca.\ ein Fünftel des lokalen Weyl-Abstands),
Vorzeichenwechsel-Detektion und Wurzelverfeinerung (\texttt{brentq} bzw.\ Illinois).
Die Entfaltung folgt dem glatten Weyl-Hauptterm
\begin{equation}
  N(t) = \frac{t}{2\pi}\,\log\!\Bigl(\frac{t}{2\pi e}\Bigr) + \frac{7}{8},
  \label{eq:weyl-unfold}
\end{equation}
und liefert $N-1$ normierte Abstände $s_n = N(\gamma_{n+1}) - N(\gamma_n)$.
Als Referenz für das zweite Moment dient $\sigma^2_{\mathrm{GUE}} = \monopolGUE$
(Wigner-Dyson-Verteilung im GUE-Limes).

%=============================================================================
\section{Numerische Ergebnisse}
\label{sec:numerische-ergebnisse}
%=============================================================================

\subsection{Stichproben $N=150$, $500$ und $1000$}

Tabelle~\ref{tab:haupt-ergebnisse} fasst drei unabhängige Läufe im Sektor
$t_0=5000$ zusammen. Der mittlere normierte Abstand bleibt in allen Fällen
nahe $1$ (Weyl-Konformität). Der Varianz-Exzess $\Delta\sigma^2 = \sigma^2 - \sigma^2_{\mathrm{GUE}}$
liegt stabil im Bereich $+0{,}053$ bis $+0{,}057$.

\begin{table}[htbp]
\centering
\caption{Spektralstatistiken nach korrigierter Pipeline ($t_0=5000$,
$\sigma^2_{\mathrm{GUE}}\approx 0{,}104$).}
\label{tab:haupt-ergebnisse}
\begin{tabular}{rcccccc}
\toprule
$N$ & $t_{\max}$ & $\bar{s}$ & $\sigma^2$ & $\Delta\sigma^2$ & $P_m$ & Laufzeit \\
\midrule
150  & \monopolTmax      & \monopolMeanSpacing  & \monopolVariance   & $+\monopolDeltaSigma$   & $-0{,}0289$ & $31{,}3\,\mathrm{s}$ \\
500  & \bulkFiveHundredTmax & \bulkFiveHundredMean & \bulkFiveHundredVar & $+\bulkFiveHundredDelta$ & $\bulkFiveHundredPm$ & $\bulkFiveHundredRuntime\,\mathrm{s}$ \\
1000 & \bulkThousandTmax    & \bulkThousandMean    & \bulkThousandVar    & $+\bulkThousandDelta$    & $\bulkThousandPm$    & $\bulkThousandRuntime\,\mathrm{s}$ \\
\bottomrule
\end{tabular}
\end{table}

\noindent
Für $N=150$ gilt $S_N=\monopolPressure$ und $P_m=S_N/150\approx -0{,}0289$.
Für $N=500$ ist $S_N=\bulkFiveHundredPressureSum$ und $P_m=\bulkFiveHundredPm$.
Die Bulk-Evaluation ($N=1000$) nutzt \texttt{mpmath\_dps=35} und Illinois-Solver
(\texttt{EABC\_Bulk\_Evaluation.py}).

\subsection{Konvergenz innerhalb eines Laufs}

Teilstichproben des $N=150$-Referenzlaufs zeigen, dass $\bar{s}$ durchgängig nahe $1$
bleibt und $\Delta\sigma^2$ sich im Bereich $+0{,}05$ bis $+0{,}07$ einpendelt
(Tabelle~\ref{tab:varianz-konvergenz}).

\begin{table}[htbp]
\centering
\caption{Konvergenz normierter Abstandsstatistiken (gleicher Lauf, $t_0=5000$).}
\label{tab:varianz-konvergenz}
\begin{tabular}{rccc}
\toprule
$N$ & $\bar{s}$ & $\sigma^2$ & $\Delta\sigma^2$ \\
\midrule
50   & $\convFiftyMean$      & $\convFiftyVar$      & $+\convFiftyDelta$      \\
100  & $\convHundredMean$     & $\convHundredVar$     & $+\convHundredDelta$     \\
150  & $\monopolMeanSpacing$  & $\monopolVariance$    & $+\monopolDeltaSigma$    \\
500  & $\bulkFiveHundredMean$  & $\bulkFiveHundredVar$  & $+\bulkFiveHundredDelta$  \\
1000 & $\bulkThousandMean$    & $\bulkThousandVar$    & $+\bulkThousandDelta$    \\
\bottomrule
\end{tabular}
\end{table}

\subsection{Block-Stabilität bei $N=1000$}

Zur Prüfung der Robustheit wurde die $N=1000$-Stichprobe in vier Blöcke à $250$
Nullstellen zerlegt (Implementierung in \texttt{EABC\_Bulk\_Evaluation.py}).
$\Delta\sigma^2$ bleibt in allen Blöcken positiv ($+0{,}049$ bis $+0{,}066$),
$P_m$ wechselt jedoch das Vorzeichen und schwankt um eine Größenordnung
--- ein Hinweis auf die oszillatorische Natur von \eqref{eq:pressure-sum},
nicht auf eine konvergente ``Monopol-Stärke''.

\begin{table}[htbp]
\centering
\caption{Blockweise Stabilität ($N=1000$, je 250 Nullstellen).}
\label{tab:block-stabilitaet}
\begin{tabular}{rcccc}
\toprule
Block & $\bar{s}$ & $\sigma^2$ & $\Delta\sigma^2$ & $P_m$ \\
\midrule
1 & $\blkOneMean$   & $\blkOneVar$   & $+\blkOneDelta$   & $\blkOnePm$   \\
2 & $\blkTwoMean$   & $\blkTwoVar$   & $+\blkTwoDelta$   & $\blkTwoPm$   \\
3 & $\blkThreeMean$ & $\blkThreeVar$ & $+\blkThreeDelta$ & $\blkThreePm$ \\
4 & $\blkFourMean$  & $\blkFourVar$  & $+\blkFourDelta$  & $\blkFourPm$  \\
\midrule
\multicolumn{5}{l}{\textit{Gesamt:} $\bar{s}=\bulkThousandMean$,
  $\sigma^2=\bulkThousandVar$,
  $\Delta\sigma^2=+\bulkThousandDelta$,
  $P_m=\bulkThousandPm$.}
\end{tabular}
\end{table}

%=============================================================================
\section{Interpretation}
\label{sec:interpretation}
%=============================================================================

\paragraph{Varianz, nicht Mittelwert.}
Die numerisch gesicherte Abweichung vom GUE-Referenzensemble manifestiert sich
im \emph{zweiten Moment} der entfalteten Abstandsverteilung:
$\Delta\sigma^2 \approx +0{,}054$--$+0{,}057$.
Der mittlere Abstand $\bar{s}\approx 1$ bleibt Weyl-konform.
Eine frühere Fehlinterpretation (Faktor $\sim 3$ im Mittelwert) resultierte
ausschließlich aus unvollständiger Siegel-Approximation
(Abschnitt~\ref{sec:methodische-richtigstellung}).

\paragraph{Relative Größe des Exzesses.}
Bei $N=1000$ beträgt $\sigma^2/\sigma^2_{\mathrm{GUE}} \approx 1{,}55$;
der Varianz-Überschuss ist damit moderat, aber über Teilstichproben und
Datenblöcke hinweg konsistent positiv. Ob dieser Exzess über größere $t$-Intervalle
persistiert oder statistisch im Rahmen endlicher GUE-Fluktuationen liegt,
bedarf weiterer Untersuchung.

\paragraph{Kein Monopol-Nachweis.}
Die Größe $P_m$ ist eine \emph{Modellgröße}, konstruiert aus den gefundenen
Nullstellen und dem Postulat \eqref{eq:qtot}. Sie belegt weder die Existenz
eines physikalischen magnetischen Monopols noch eine kausale Kopplung zwischen
Hurwitz-Gitter-Defekten und Riemann-Nullstellen. Korrelationen zwischen
$\Delta\sigma^2$ und $P_m$ wurden in diesem Preprint \emph{nicht} quantitativ
getestet.

\paragraph{Zurückhaltung bei Mechanismen.}
Zuordnungen von $\Delta\sigma^2$ zu konkreten Feldtheorie-Bildern
(Stark-Effekt, Symmetriebrechung der Dichte, Poisson-zu-GUE-Übergang mit
zusätzlicher Breite) bleiben spekulativ und werden bis zu unabhängigen
Reproduktionen zurückgestellt.

%=============================================================================
\section{Offene Fragen}
\label{sec:offene-fragen}
%=============================================================================

\begin{enumerate}
  \item \textbf{Herleitung von $Q_{\mathrm{tot}}=1796$.}
        Welche arithmetische oder physikalische Begründung (jenseits des Modellpostulats)
        rechtfertigt den Wert $1796$? Eine Verbindung zu E$_8$ ist nicht etabliert.

  \item \textbf{Vereinheitlichung der Monopol-Bilder.}
        Wie lässt sich das globale Punkt-Monopol mit den Hurwitz-Defekten aus
        \texttt{MagMo.py} in einem gemeinsamen Kopplungsterm formulieren?

  \item \textbf{Dirac-Monopol und Spektralstatistik.}
        Unter welchen Annahmen würde ein Dirac-monopolarer Störterm die
        \emph{Varianz} (nicht die mittlere Dichte) der Nullstellen-Abstände
        erhöhen, ohne das Weyl-Gesetz zu verletzen?

  \item \textbf{t'Hooft-Skala und $Q_{\mathrm{tot}}$.}
        Lässt sich die t'Hooft-Monopolmasse $M_M\sim M_X/\alpha$
        (Abschnitt~\ref{sec:hooft-monopolmasse}) mit dem Postulat
        $Q_{\mathrm{tot}}=1796$ oder mit $\alpha_G^{\mathrm{EABC}}$ in einer
        gemeinsamen EABC-Kopplungsgleichung formulieren?

  \item \textbf{Skalenabhängigkeit.}
        Bleibt $\Delta\sigma^2 > 0$ für größere $t$ und $N$ stabil?
        Systematische Fenster-Scans oberhalb $t\approx 6000$ fehlen.

  \item \textbf{Unabhängige Reproduktion.}
        Die Pipeline ist in Python/mpmath dokumentiert; externe Replikation
        mit alternativem Nullstellen-Solver (Odlyzko-Schönage, arb) wäre wünschenswert.
\end{enumerate}

%=============================================================================
\section*{Danksagung und Hinweise}
\addcontentsline{toc}{section}{Danksagung und Hinweise}
%=============================================================================

Die numerischen Skripte \texttt{EABC\_Gross\_Analyse.py},
\texttt{EABC\_Bulk\_Evaluation.py}, \texttt{Monopol Test.py}, \texttt{MagMo.py}
und \texttt{hooft.py} sind die primären Reproduktionsquellen. Ergebnismakros für $N=150$ stammen aus
\texttt{monopol\_preprint\_results.txt}; Bulk-Daten aus \texttt{bulk\_results.json}.

%=============================================================================
\begin{thebibliography}{9}
\bibitem{berry-keating}
J.~P.~Keating und M.~C.~Gutzwiller,
Spektralinterpretationen der Riemann-Nullstellen (Überblick).

\bibitem{gue}
M.~L.~Mehta,
\emph{Random Matrices},
Academic Press, 3.\ Aufl.

\bibitem{odlyzko}
A.~M.~Odlyzko,
On the distribution of spacings between zeros of the zeta function,
\emph{Math.\ Comp.} \textbf{48} (1987).

\bibitem{dirac}
P.~A.~M.~Dirac,
Quantised singularities in the electromagnetic field,
\emph{Proc.\ Roy.\ Soc.\ A} \textbf{133} (1931).

\bibitem{thooft}
G.~'t~Hooft,
Magnetic monopoles in unified gauge theories,
\emph{Nucl.\ Phys.\ B} \textbf{79} (1974), 276--284.

\bibitem{polyakov}
A.~M.~Polyakov,
Particle spectrum in quantum field theory,
\emph{JETP Lett.} \textbf{20} (1974), 194--195.

\bibitem{hurwitz}
A.~Hurwitz,
Über die Komposition der quadratischen Formen,
\emph{Math.\ Ann.} \textbf{88} (1923).
\end{thebibliography}

\end{document}
